% SEGMENT OLP-0548-S01
% Part: set-theory
% Chapter: ordinals
% Section: introduction

\documentclass[../../../include/open-logic-section]{subfiles}

\begin{document}

\olfileid{sth}{ordinals}{intro}

\olsection{அறிமுகம்}

\olref[z][]{chap} இல், படிநிலை வரிசையில் ஒரு முடிவிலிப் படிநிலை உள்ளது
என்று \stagesinf{} வடிவில் நாம் எடுகோளிட்டோம் (நமது முடிவிலி அடிகோளையும்
பார்க்கவும்). ஆனால் \stagessucc{} என்பதை ஏற்றுக்கொண்டால், முடிவிலிப்
படிநிலையோடு நாம் நின்றுவிட முடியாது; தொடர்ந்து செல்ல வேண்டும். ஆகவே,
முதல் முடிவிலிப் படிநிலைக்கு அடுத்த படிநிலையில், அந்த முதல் முடிவிலிப்
படிநிலையில் கிடைத்த கணங்களைக் கொண்டு அமைக்கக்கூடிய எல்லாத் தொகுப்புகளையும்
உருவாக்குகிறோம்; மீண்டும் இதைச் செய்கிறோம்; மீண்டும்; மீண்டும்; \dots

இங்கு மறைமுகமாக நிகழ்ந்திருப்பது இதுதான்: எல்லா இயல் எண்களுக்கும்
\emph{பின்னால்} எண்கள் இருக்கலாம் என்ற ஓர் ``உள்ளுணர்வு'' எண் கருத்தை நாம்
பயன்படுத்தத் தொடங்கியுள்ளோம். குறிப்பாக, இங்கு பயன்படும் கருத்து
\emph{முடிவிலிக்கடந்த வரிசையெண்} என்பதாகும். இந்தக் கருத்தை மேலும்
துல்லியமாக்குவதே இவ்வத்தியாயத்தின் நோக்கம். வரிசையெண்ணின் பொதுக் கருத்தை
ஆராய்ந்த பின், சில குறிப்பிட்ட கணங்களையே நமது வரிசையெண்களாக வெளிப்படையாக
வரையறுப்போம்.

\end{document}
